% ============================================================================
\clearpage
\section*{v4.10.0-r07 Source-Only: The Relative-Entropy Chain in the Tachyon Frame, and a Correction to r06.5}
\addcontentsline{toc}{section}{v4.10.0-r07 Source-Only: The Relative-Entropy Chain in the Tachyon Frame, and a Correction to r06.5}
\label{sec:v4100-r07-tachyon-chain}
% ============================================================================

\paragraph{Purpose.}
The author asked whether the companion can map the derivation of Dorau and Much (\S\ref{sec:v4100-r06-relative-entropy}) onto its own frames and
so obtain the Einstein equations internally.  A pre-registered test was planned.  The search of prior work showed that the answer is fixed
analytically, so this layer is an exploration, not a test; no hypothesis number is assigned.  On the way, Remark r06.5 turned out to be wrong.  A second question of the author, why frame independence is reached and the prefactor of the area law is not, is answered in r07.4.
The notes are \texttt{notes/\allowbreak research/\allowbreak v4\_10\_0\_r07/}.  Earlier layers are not rewritten.

\noindent\textbf{Source-only status.}
This v4.10.0-r07 layer is source-only.  It does not build a PDF, does not assign a new DOI,
and does not authorize a public release.  The v4.10.0 Version DOI
\href{https://doi.org/10.5281/zenodo.23991357}{10.5281/zenodo.23991357} and the Concept DOI
\href{https://doi.org/10.5281/zenodo.19210728}{10.5281/zenodo.19210728} are unchanged by it.

\subsection*{v4.10.0-r07 results}
\label{subsec:v4100-r07-results}

\paragraph{Correction v4.10.0-r07.1 (to Remark r06.5).}
Remark v4.10.0-r06.5 states that under the tachyon map \(z\mapsto iz\) the rapidity becomes imaginary and the boost a rotation, and asks whether
\(z\mapsto iz\) is the Wick rotation of Theorem~\ref{thm:artqft}.  This is wrong.  With \(z=E+iB\) (v4.9.0-r02), \(z\mapsto iz\) is the
electromagnetic duality rotation \(E\mapsto-B\), \(B\mapsto E\), a map of field space, not of spacetime.  By Theoremlet v4.8.0-r20.1 it commutes with
every element of \(SO(3,\mathbb C)\), the proper Lorentz group acting on \(z\).  A boost stays a boost with real rapidity.  The question of r06.5 is
answered in the negative: \(z\mapsto iz\) is not the Wick rotation \(t\mapsto i\tau_E\); the two act on different objects.  The frame entries of r06 for
frame T are replaced by r07.2.

\paragraph{Theoremlet v4.10.0-r07.2 (the chain in frame T for the free Maxwell field).}
Exploration, not blind (fixed analytically before the check).  For the free Maxwell field, the chain of r06.1 gives in frame T the same Einstein
equation as in frame M: the same coupling, the same sign, the same inputs.
\begin{enumerate}[leftmargin=2em]
\item[(a)] The duality rotation \(U_\alpha\) is a symmetry of the free Maxwell field (conserved helicity).  It acts pointwise on \(F\), maps the algebra of
the Rindler wedge onto itself and leaves the vacuum invariant.
\item[(b)] It commutes with the boost (v4.8.0-r20.1), hence with the modular group of the vacuum (Bisognano--Wichmann).
\item[(c)] Relative entropy is invariant under such an automorphism: \(S_{\rm rel}(\omega_0\Vert\omega_\phi)=S_{\rm rel}(\omega_0\Vert\omega_{U\phi})\).
\item[(d)] The energy--momentum tensor is Hermitian in \(z\) (energy density \(\tfrac12|z|^2\), Poynting vector \((\bar z\times z)/2i\)) and therefore
duality invariant; the flux \(T_{ab}\xi^a\xi^b\) is unchanged.  The bilinear invariant \(s=z\cdot z\) changes sign: the Lagrangian density
\(\tfrac12(E^2-B^2)\) and \(E\cdot B\) are odd (v4.8.0-r06.3, v4.9.0-r40).
\end{enumerate}
\emph{Reading.} Gravity couples to the Hermitian tensor; the frames differ in the bilinear invariants.  For the Maxwell field, gravity is therefore
frame blind.  This turns the reading of v4.9.0-r64 (``gravity itself is frame independent'') into a statement for this case.  It is not new physics:
the duality invariance of the Maxwell tensor is textbook knowledge.

\paragraph{Proof and measurement sources.}
(a)--(c) are standard (duality symmetry of the free Maxwell field; invariance of the Araki relative entropy under automorphisms that fix the
reference state).  (d) and the commutation are checked numerically by \texttt{k217\_\allowbreak dualitaet\_\allowbreak gravitation.py} (2000 random fields and
boosts): commutation, invariance of the tensor and of the null flux, and the sign change of \(E^2-B^2\) and \(E\cdot B\) hold exactly in floating point;
a non-unitary scaling \((1+0.3i)z\) changes the tensor by \(O(1)\) (negative control).  \(\square\)

\paragraph{Theoremlet v4.10.0-r07.3 (bookkeeping: what the internal version gains).}
The chain is written in QAM bookkeeping for frames M and T, with the duality rotation as a declared map
(\texttt{qam\_\allowbreak dorau\_\allowbreak much\_\allowbreak companion.tex}); \texttt{qam.sty} raises no error.  The exported record gives for both frames the
Einstein equation with status \emph{conditional} and inputs \{local flatness, area law\}, the same inputs as in Dorau--Much.  The internal version
gains frame independence as a statement (r07.2).  It does not reduce the inputs: the area law, and with it the normalization \(1/4G\), remains an
input (r06.2).  For comparison, the same bookkeeping applied to the paper itself (\texttt{qam\_\allowbreak dorau\_\allowbreak much\_\allowbreak paper.tex}) reports a
promotion if the Einstein equation is declared proved, the area law as an input of the result, and, under the main-body typing of the
Bekenstein--Hawking law as a photon-frame statement (\S\ref{sec:bh}), a frame mixing without declared map.

\paragraph{Theoremlet v4.10.0-r07.4 (why frame independence is reached and the prefactor of the area law is not).}
Analytic, not blind (author question after r07.3).
\begin{enumerate}[leftmargin=2em]
\item[(a)] \emph{Frame independence is a statement about relations.}  The map between M and T is a symmetry of the system that fixes the reference state
and commutes with its modular group (r07.2), and every link of the chain is built from objects invariant under it.  Such an argument yields equalities
between the frames.  It says nothing about the value both frames share.
\item[(b)] \emph{The prefactor is an absolute number.}  The area law \(S=\delta A/4\) fixes the entropy per Planck area, i.e.\ the coupling (r06.2).  A
symmetry between frames cannot fix it; something must give one frame its absolute value.
\item[(c)] \emph{The countings used by the companion are calibrated to it.}  A horizon counting gives \(S=(\kappa/\gamma)\,A/\ell_P^2\) with a combinatorial
constant \(\kappa\); the value of the Barbero--Immirzi parameter at which it reproduces the area law \(S=A/(N\ell_P^2)\) is \(\gamma_*=N\kappa\).  Both literature
values used by the companion are of this kind: \(\gamma_{\rm string}=4\kappa_{\rm ABCK}\), \(\kappa_{\rm ABCK}=\ln2/(4\pi\sqrt3)\) (v4.8.0-r03 F11a), and
\(\gamma_{\rm BI}=4\kappa_{\rm DL/M}\).  Their ratio
\[
\gamma_L=\frac{\gamma_{\rm BI}}{\gamma_{\rm string}}=\frac{\kappa_{\rm DL/M}}{\kappa_{\rm ABCK}}=1.864700
\]
does not depend on \(N\): the prefactor cancels (checked for \(N=1,4,8\)).  So do \(\beta_{\rm sym}=(\gamma_L-1)/(\gamma_L+1)\) and \(\Omega_M=\beta_{\rm sym}^2\)
(H96).  Every quantity the companion builds from the two countings is exactly the part in which the normalization of the area law drops out.
\item[(d)] \emph{Consequence.}  With these inputs the prefactor cannot be derived, in principle and not only in practice: it was put into each
\(\gamma_*\) and removed again by taking the ratio.  The companion inherits the Barbero--Immirzi ambiguity of loop quantum gravity, in which the area
law is matched, not derived.  A derivation would need one absolute value of \(\gamma\) (or of the entropy per state) from the structure of the companion,
not a ratio; no such statement exists in the corpus.
\item[(e)] \emph{A caution.}  The state disc at the Hopf point has eigenvalues \(3/4\) and \(1/4\) (v4.9.0-r17).  Identifying that \(1/4\) with the one of the
area law would be a single-number identification of the kind recorded in v4.10.0-r05.2; there is no mechanism linking the two, and it is not proposed.
\end{enumerate}

\paragraph{Limits.}
Only the free Maxwell field.  For massive fields the tachyon image of the action is tachyonic (H71); a tachyonic field has no stable Hadamard
ground state and the premises of the chain fail.  Under the second reading (v4.9.0-r84, H79) the tachyon part is a weight on the matter action; how it
enters the flux is not examined.  Interacting fields: duality is in general not a symmetry.

\paragraph{Consequence.}
The companion reproduces the Dorau--Much derivation in both frames for the Maxwell field and so obtains frame-blind gravity as a statement.  It adds
no derivation of the area law; \(G\) and \(\Lambda\) remain as in r06.  The reason is structural (r07.4): the frame map gives relations,
and the countings enter only through a ratio in which the prefactor cancels.  The open front is one absolute value of \(\gamma\) from the structure.

\subsection*{v4.10.0-r07 status in the three frames}
\label{subsec:v4100-r07-frames}

\begin{itemize}[leftmargin=2em]
\item M: the chain of r06.1 holds as stated.
\item T: for the free Maxwell field the chain holds unchanged (r07.2); for massive fields the tachyon image has no stable ground state (open under
the second reading).
\item P: the null sector \(s=0\) is fixed by \(z\mapsto iz\); no separate statement.
\end{itemize}
For the Maxwell field, M and T agree: the Einstein equation of the chain is frame free in that case.

\subsection*{v4.10.0-r07 register entries}
\label{subsec:v4100-r07-register}

\begin{center}
\small
\begin{tabular}{@{}>{\raggedright\arraybackslash}p{0.34\textwidth}>{\raggedright\arraybackslash}p{0.58\textwidth}@{}}
\toprule
statement & status after v4.10.0-r07\\
\midrule
r06.5 (boost becomes rotation under \(z\mapsto iz\)) & wrong; \(z\mapsto iz\) commutes with \(SO(3,\mathbb C)\) (correction r07.1)\\
Einstein equation in frame T (Maxwell) & same as in frame M (exploration, not blind)\\
gravity frame independent (r64) & a statement for the free Maxwell field; open for massive fields\\
inputs of the internal chain & same as Dorau--Much (local flatness, area law)\\
prefactor of the area law & not derivable from the two countings: it cancels in \(\gamma_L\) (r07.4)\\
\bottomrule
\end{tabular}
\end{center}

\subsection*{v4.10.0-r07 ledger and guard update}
\label{subsec:v4100-r07-ledger-guard}

\[
R_{\rm chk}^{v4100\text{-}r07}=(1,1,0)
\]
for (r06.5 corrected; chain in frame T for the Maxwell field; area law still an input, prefactor cancels in the ratio of countings).

The layer spends no guard:
\[
R_{\rm glob}^{(v4100\text{-}r07,+)}=(0,1,1,1,1),
\qquad
\Gamma_{\rm guard}^{v4100\text{-}r07}=(1,1,1,1,1),
\qquad
\beta_4=1/256\ \text{unspent}.
\]
It does not imply exact finite \(C5\) zero closure, exact all-readout terminality,
\(C6/T2\), a global photon-clock/RH theorem, or global all-readout closure.
